%! Composition of Functions — Unit 1, Lesson 1.5 — Student Packet Cover Sheet
\documentclass[10pt]{article}
\usepackage{precalculus-article}
\usepackage{precalculus-boxes}
\usepackage{ltablex}
\keepXColumns

\begin{document}

% Full-bleed plum banner
\begin{tikzpicture}[remember picture, overlay]
  \fill[plum] (current page.north west)
    rectangle ([yshift=-0.9in]current page.north east);
\end{tikzpicture}
\vspace{-0.8in}

\begin{center}
  {\color{white}\textbf{\LARGE Precalculus}}\\[4pt]
  {\color{lilacmid}\large\textbf{Unit 1: Functions and Change}}\\[6pt]
  {\color{white}\textbf{\large Lesson 1.5 \quad Composition of Functions}}
\end{center}

\vspace{0.22in}
\namedateperiod
\vspace{0.18in}

\begin{learningtargetbox}
\small
\begin{enumerate}[leftmargin=1.5em, itemsep=5pt]
  \item I can find the \textit{handoff} in a composite --- the number the inner
        function gives to the outer one --- and use it to evaluate $f(g(3))$
        from a table or from formulas.
  \item I can build $(f \circ g)(x)$ by substituting $g(x)$ for every $x$ in
        $f(x)$, and show that swapping the order usually changes the function.
  \item I can explain why a composite has no value when the handoff is not
        something the outer function accepts.
  \item I can decompose a function into an inner and an outer piece, and check
        my split by substituting back.
  \item I can rewrite Lesson 1.4's ``inside'' changes as compositions, and say
        why that is what makes them act on inputs.
\end{enumerate}
\end{learningtargetbox}

\vspace{0.14in}

\begin{tocbox}
\small
\renewcommand{\arraystretch}{1.5}
\begin{tabularx}{\linewidth}{@{} c l X r @{}}
\rowcolor{lilac}
\textbf{\#} & \textbf{Component} & \textbf{Description} & \textbf{Score} \\
\midrule
1 & Warm-Up        & Spiral review of prerequisite skills              & \blank{1.2cm} \\
2 & Guided Notes   & The handoff: find it, hand it over, take it apart & \blank{1.2cm} \\
3 & Group Activity & The summit cable car: minutes in, degrees out     & \blank{1.2cm} \\
4 & Exit Ticket    & Independent check on evaluating, building, decomposing & \blank{1.2cm} \\
5 & Homework       & Practice and extension                            & \blank{1.2cm} \\
\midrule
  & \multicolumn{2}{r}{\textbf{Total}} & \blank{1.2cm} \\
\end{tabularx}
\end{tocbox}

\vspace{0.14in}

\begin{remindbox}
\small
Last lesson a rule interfered with the \textit{question} you hand a function or
the \textit{answer} it hands back. Today two functions work in a row, and one
number sits between them: the \textbf{handoff}. It is the inner function's
answer and the outer function's question at the same time --- so find it first
and write it down. Everything else follows from it. Swap the machines and you
swap the handoff, which is why $f(g(x))$ and $g(f(x))$ are usually different
functions. And a handoff the outer function will not accept means the composite
has no value at all.
\end{remindbox}

\end{document}
